% practice.tex -- Module 5 (Advanced Mechanics): Practice set.
% 24 new HSC-style questions: Section I 8 multiple choice, Section II 10 short answer,
% Section III 6 extended response. 74 marks. Ordered from easy to HSC-hard.
% Answers: practice_answers.tex (one practiceans per practiceq, same order).
% All questions are new (new contexts and numbers); none is copied from a booklet,
% homework or mock paper. Q23 uses published TRAPPIST-1 orbital data (semi-major axes and
% periods, Agol et al. 2021), rounded to 3 significant figures.
% \Needspace before each multiple-choice question keeps its stem and options together
% (conditional, layout-independent; no \clearpage or \pagebreak). Q13 and Q16 (side-by-side
% diagrams) and the Section III heading also have one, so no header is left alone at a page foot.
% Local helpers (prefix "prac"), defined below: \pracmarks, \pracblank.
% Figures and tables here use \caption* (no number): this chapter is unnumbered, so a
% numbered caption would continue the numbering of Chapter 5.
\providecommand\pracmarks[1]{\hfill\mbox{\textcolor{hscmuted}{\small(#1)}}}
% \pracblank: an empty cell for the student to fill in (a short grey line).
\providecommand\pracblank{\textcolor{hscmuted}{\rule[-2pt]{2.6em}{0.4pt}}}

\unnumberedchapter[\faDumbbell]{Practice set}

These 24 questions are written in the style of the HSC exam. They test all five chapters of
this book. They start easy and get harder.
\begin{itemize}
  \item \textbf{Section I:} 8 multiple-choice questions, 1 mark each. Choose A, B, C or D.
  \item \textbf{Section II:} 10 short-answer questions, 3 or 4 marks each.
  \item \textbf{Section III:} 6 extended-response questions, 4 to 7 marks each.
\end{itemize}
The total is 74 marks. Allow about 2 hours and 15 minutes. (In the HSC exam you have about
1.8 minutes for each mark.)

\textbf{For every question:}
\begin{itemize}
  \item Use the formula sheet and these data-sheet values: $g = \qty{9.8}{m.s^{-2}}$,
        $G = \qty{6.67e-11}{N.m^2.kg^{-2}}$, $M_\mathrm{E} = \qty{6.0e24}{kg}$,
        $r_\mathrm{E} = \qty{6.371e6}{m}$.
  \item Ignore air resistance, unless the question tells you otherwise.
  \item Show your working. Give each answer with a unit and sensible significant figures.
\end{itemize}

The marks for each part of a question are shown on the right, like this: \textcolor{hscmuted}{\small(2 marks)}.
The tag next to each question number (for example \textcolor{hscbrand}{Ch~2}) shows which chapter
the question tests. If you get a question wrong, revise that chapter. Full worked answers and
marking criteria are at the back of the book. Click a question number to jump to its answer.

% =====================================================================
\section*{Section I: Multiple choice}

\Needspace{11\baselineskip}
\begin{practiceq}{1}{1}\label{pq:top}
A ball is thrown at an angle above the horizontal. At the highest point of its path, which
statement is correct?
\choices[1]{Its velocity is zero, and its acceleration is zero.}
  {Its velocity is zero, and its acceleration is \qty{9.8}{m.s^{-2}} downwards.}
  {Its velocity is horizontal, and its acceleration is zero.}
  {Its velocity is horizontal, and its acceleration is \qty{9.8}{m.s^{-2}} downwards.}
\end{practiceq}

\Needspace{18\baselineskip}
\begin{practiceq}{1}{1}\label{pq:vygraph}
A ball is launched from level ground at an angle above the horizontal. It lands at the same
height. Up is positive. Which graph shows the vertical velocity $v_y$ of the ball against
time $t$, from launch until it lands?
\choices[2]{%
  \begin{tikzpicture}[x=1cm, y=0.5cm, baseline=(current bounding box.north)]
    \draw[axisline] (0,0) -- (2.8,0) node[right, font=\small] {$t$};
    \draw[axisline] (0,-1.6) -- (0,1.7) node[above, font=\small] {$v_y$};
    \node[left, font=\footnotesize, text=hscmuted] at (0,0) {0};
    \draw[hscseriesA, line width=1.3pt] (0,1.2) -- (2.4,-1.2);
  \end{tikzpicture}}
  {%
  \begin{tikzpicture}[x=1cm, y=0.5cm, baseline=(current bounding box.north)]
    \draw[axisline] (0,0) -- (2.8,0) node[right, font=\small] {$t$};
    \draw[axisline] (0,-1.6) -- (0,1.7) node[above, font=\small] {$v_y$};
    \node[left, font=\footnotesize, text=hscmuted] at (0,0) {0};
    \draw[hscseriesA, line width=1.3pt] (0,1.2) -- (1.2,0) -- (2.4,1.2);
  \end{tikzpicture}}
  {%
  \begin{tikzpicture}[x=1cm, y=0.5cm, baseline=(current bounding box.north)]
    \draw[axisline] (0,0) -- (2.8,0) node[right, font=\small] {$t$};
    \draw[axisline] (0,-1.6) -- (0,1.7) node[above, font=\small] {$v_y$};
    \node[left, font=\footnotesize, text=hscmuted] at (0,0) {0};
    \draw[hscseriesA, line width=1.3pt] (0,1.2) -- (2.4,1.2);
  \end{tikzpicture}}
  {%
  \begin{tikzpicture}[x=1cm, y=0.5cm, baseline=(current bounding box.north)]
    \draw[axisline] (0,0) -- (2.8,0) node[right, font=\small] {$t$};
    \draw[axisline] (0,-1.6) -- (0,1.7) node[above, font=\small] {$v_y$};
    \node[left, font=\footnotesize, text=hscmuted] at (0,0) {0};
    \draw[hscseriesA, line width=1.3pt, domain=0:2.4, samples=41, smooth]
      plot (\x, {4*1.3*(\x/2.4)*(1-\x/2.4)});
  \end{tikzpicture}}
\end{practiceq}

\Needspace{8\baselineskip}
\begin{practiceq}{1}{2}\label{pq:ratio}
A car goes around a flat bend at a constant speed. Friction provides its centripetal force,
$F$. A second car has the same mass. It goes around a bend with \emph{half} the radius, at
\emph{twice} the speed. What centripetal force does the second car need?
\choices{$2F$}{$4F$}{$8F$}{$16F$}
\end{practiceq}

\Needspace{11\baselineskip}
\begin{practiceq}{1}{2}\label{pq:banked}
A car goes around a banked curve at exactly its design speed, so it needs no sideways
friction. Which force provides the centripetal force on the car?
\choices[2]{the weight of the car}{the friction force between the tyres and the road}
  {the vertical component of the normal force}{the horizontal component of the normal force}
\end{practiceq}

\Needspace{11\baselineskip}
\begin{practiceq}{1}{3}\label{pq:merry}
A merry-go-round (a round playground ride that spins) turns at a steady rate. Child P sits
\qty{1.0}{m} from the centre. Child Q
sits \qty{2.0}{m} from the centre. Which statement is correct?
\choices[1]{P and Q have the same speed.}{Q has twice the angular velocity of P.}
  {Q has twice the centripetal acceleration of P.}
  {Q has four times the centripetal acceleration of P.}
\end{practiceq}

\Needspace{12\baselineskip}
\begin{practiceq}{1}{3}\label{pq:door}
A door is hinged on one side. Four people push on it, one at a time. Each angle is measured
between the force and the door. Which push gives the largest torque about the hinge?
\choices[1]{\qty{30}{N}, \qty{0.90}{m} from the hinge, at \ang{90} to the door}
  {\qty{50}{N}, \qty{0.60}{m} from the hinge, at \ang{90} to the door}
  {\qty{80}{N}, \qty{0.90}{m} from the hinge, at \ang{20} to the door}
  {\qty{100}{N}, \qty{0.25}{m} from the hinge, at \ang{90} to the door}
\end{practiceq}

\Needspace{8\baselineskip}
\begin{practiceq}{1}{4}\label{pq:planetg}
A planet has 4 times the mass of the Earth and 2 times the radius of the Earth. At the
Earth's surface, $g = \qty{9.8}{N.kg^{-1}}$. What is the gravitational field strength at the
surface of the planet?
\choices{\qty{2.45}{N.kg^{-1}}}{\qty{9.8}{N.kg^{-1}}}{\qty{19.6}{N.kg^{-1}}}{\qty{39.2}{N.kg^{-1}}}
\end{practiceq}

\Needspace{12\baselineskip}
\begin{practiceq}{1}{5}\label{pq:higher}
A satellite moves from one circular orbit around the Earth to a \emph{higher} circular orbit.
How do its kinetic energy $K$, its gravitational potential energy $U$ and its total energy
$E$ change?
\choices[1]{$K$ decreases, $U$ increases and $E$ increases.}
  {$K$ decreases, $U$ increases and $E$ stays the same.}
  {$K$ increases, $U$ decreases and $E$ decreases.}
  {$K$ stays the same, $U$ increases and $E$ increases.}
\end{practiceq}

% =====================================================================
\section*{Section II: Short answer}

\begin{practiceq}{3}{1}\label{pq:drone}
A delivery drone flies horizontally at a steady \qty{7.5}{m.s^{-1}}, \qty{20}{m} above flat
ground. It releases a small first-aid kit.
\begin{enumerate}[label=(\alph*)]
  \item Calculate the time the kit takes to reach the ground. \pracmarks{1 mark}
  \item Calculate the horizontal distance the kit travels before it lands. \pracmarks{1 mark}
  \item The drone keeps flying at the same velocity. Where is the drone when the kit lands?
        Explain your answer. \pracmarks{1 mark}
\end{enumerate}
\end{practiceq}

\begin{practiceq}{3}{2}\label{pq:centrifuge}
Pilots train in a centrifuge: a seat on the end of a long arm that spins in a horizontal
circle. The radius of the circle is \qty{8.0}{m}. A \qty{72}{kg} pilot has a centripetal
acceleration of $6.0g$, that is $6.0 \times \qty{9.8}{m.s^{-2}}$.
\begin{enumerate}[label=(\alph*)]
  \item Calculate the speed of the pilot. \pracmarks{1 mark}
  \item Calculate the period of the motion. \pracmarks{1 mark}
  \item Calculate the size of the centripetal force on the pilot. \pracmarks{1 mark}
\end{enumerate}
\end{practiceq}

% layout QA: keep this short 3-part question on one page
\Needspace{13\baselineskip}
\begin{practiceq}{3}{3}\label{pq:turbine}
The tip of each blade of a wind turbine is \qty{45.0}{m} from the axis. The blades turn at a
steady 16.0 revolutions per minute (rpm).
\begin{enumerate}[label=(\alph*)]
  \item Calculate the angular velocity of the blades in \si{rad.s^{-1}}. \pracmarks{1 mark}
  \item Calculate the speed of a blade tip in \si{km.h^{-1}}. \pracmarks{1 mark}
  \item Point M is halfway along a blade. Compare the angular velocity and the speed of M with
        those of the tip. \pracmarks{1 mark}
\end{enumerate}
\end{practiceq}

\begin{practiceq}{3}{4}\label{pq:period}
A weather satellite moves in a circular orbit around the Earth. Its period is
\qty{100.0}{min}. Calculate its altitude (its height above the Earth's surface) in
kilometres.
\end{practiceq}

\Needspace{14\baselineskip}
\begin{practiceq}{3}{3}\label{pq:forearm}
\begin{sidebyside}[0.40]
  \centering
  \begin{tikzpicture}[x=1cm, y=1cm]
    % upper arm (vertical) and forearm (horizontal); elbow = pivot at (0,0)
    \fill[hscmuted!25, draw=hscink, line width=0.6pt, rounded corners=2pt]
      (-0.22,0.18) rectangle (0.22,2.4);
    \fill[hscmuted!25, draw=hscink, line width=0.6pt, rounded corners=2pt]
      (-0.15,-0.18) rectangle (4.25,0.18);
    \fill[hscink] (0,0) circle (2pt);
    \node[diagramlabel, below, inner sep=4pt, font=\footnotesize] at (0,-0.18) {elbow (pivot)};
    % biceps force at 0.040 m (drawn at 0.4 cm for clarity, not to scale)
    \draw[force] (0.45,0.18) -- (0.45,1.75) node[above right, inner sep=1pt, text=hscforce]
      {$F_\mathrm{b}$};
    % forearm weight at 0.16 m
    \draw[force] (1.84,0) -- (1.84,-1.0) node[below, inner sep=2pt, text=hscforce, font=\small]
      {forearm};
    % dumbbell at 0.34 m
    \fill[hscink!80] (3.85,-0.42) rectangle (3.95,0.42);
    \fill[hscink!80] (3.7,-0.3) rectangle (4.1,0.3);
    \draw[force] (3.9,-0.42) -- (3.9,-1.4) node[below, inner sep=2pt, text=hscforce, font=\small]
      {dumbbell};
  \end{tikzpicture}
  \captionof*{figure}{An arm holding a dumbbell. The elbow is the pivot.}
\nextside
A student holds a \qty{6.0}{kg} dumbbell in her hand. Her forearm is horizontal and does not
move (see the diagram).
\begin{itemize}
  \item The dumbbell is \qty{0.34}{m} from the elbow.
  \item The forearm and hand have a mass of \qty{1.8}{kg}. Their weight acts
        \qty{0.16}{m} from the elbow.
  \item The biceps muscle pulls straight up on the forearm, \qty{0.040}{m} from the elbow.
\end{itemize}
Calculate the force $F_\mathrm{b}$ that the biceps muscle exerts.
\end{sidebyside}
\end{practiceq}

\begin{practiceq}{4}{1}\label{pq:golf}
A golf ball is hit from level ground at \qty{126}{km.h^{-1}}, at \ang{24} above the
horizontal.
\begin{enumerate}[label=(\alph*)]
  \item Calculate the maximum height of the ball. \pracmarks{2 marks}
  \item Calculate the range of the ball. \pracmarks{2 marks}
\end{enumerate}
\end{practiceq}

\begin{practiceq}{3}{5}\label{pq:mercury}
Mercury has a mass of \qty{3.30e23}{kg} and a radius of \qty{2.44e6}{m}.
\begin{enumerate}[label=(\alph*)]
  \item Calculate the escape velocity from the surface of Mercury. \pracmarks{2 marks}
  \item A spacecraft is in a circular orbit just above Mercury's surface. Show that it must
        increase its speed by a factor of $\sqrt{2}$ to escape. \pracmarks{1 mark}
\end{enumerate}
\end{practiceq}

\Needspace{16\baselineskip}
\begin{practiceq}{4}{2}\label{pq:swingride}
\begin{sidebyside}[0.42]
  \centering
  \begin{tikzpicture}[x=0.5cm, y=0.5cm]
    % scale: 1 m = 0.5 cm. d = 4.0 m, L = 5.0 m at 40 deg to the vertical
    \fill[ground] (-1.6,-0.3) rectangle (9.4,0);
    \draw[hscink, line width=0.6pt] (-1.6,0) -- (9.4,0);
    \draw[guide] (0,-0.3) -- (0,8.6) node[above, diagramlabel, font=\footnotesize] {axis};
    \fill[hscmuted!35, draw=hscink, line width=0.6pt] (-0.25,0) rectangle (0.25,7.6);
    \draw[hscink, line width=1.6pt] (0,7.6) -- (4.0,7.6);
    \fill[hscink] (4.0,7.6) circle (2pt);
    \draw[hscink, line width=0.8pt] (4.0,7.6) -- (7.214,3.770);
    \draw[guide] (4.0,7.6) -- (4.0,3.0);
    \draw[hscink, line width=0.6pt] (4.0,6.0) arc (-90:-50:1.6);
    \node[diagramlabel, font=\small] at (4.55,5.35) {$\theta$};
    \draw[object] (6.85,3.25) rectangle (7.6,3.77);
    \draw[displacement, <->] (0,8.3) -- (4.0,8.3) node[midway, above, text=hscink, font=\small] {$d$};
    \node[diagramlabel, font=\small, anchor=south west] at (5.75,5.9) {$L$};
    \draw[hscink, <->, line width=0.6pt] (0,2.4) -- (7.214,2.4)
      node[midway, below, font=\small] {$r$};
    \draw[guide] (7.214,3.2) -- (7.214,2.2);
  \end{tikzpicture}
  \captionof*{figure}{A swing ride (side view, to scale).}
\nextside
In a swing ride, each chair hangs from a chain fixed to the end of a horizontal arm
(see the diagram). The arm turns around a vertical axis. The arm length is
$d = \qty{4.0}{m}$ and the chain length is $L = \qty{5.0}{m}$. When the ride turns at a
steady rate, each chain makes an angle $\theta = \ang{40}$ with the vertical. Treat each rider
and chair as one small object.
\begin{enumerate}[label=(\alph*)]
  \item Calculate $r$, the radius of the circle in which a rider moves. \pracmarks{1 mark}
  \item Calculate the speed of a rider. \pracmarks{2 marks}
  \item Calculate the period of the ride. \pracmarks{1 mark}
\end{enumerate}
\end{sidebyside}
\end{practiceq}

\begin{practiceq}{4}{3}\label{pq:vertical}
A student swings a \qty{0.20}{kg} ball on a light string in a vertical circle of radius
\qty{0.80}{m}. At the top of the circle the ball moves at \qty{3.5}{m.s^{-1}}.
\begin{enumerate}[label=(\alph*)]
  \item Calculate the tension in the string at the top of the circle. \pracmarks{2 marks}
  \item Use energy to calculate the speed of the ball at the bottom of the circle. Assume
        that only gravity does work on the ball. \pracmarks{1 mark}
  \item Calculate the tension in the string at the bottom of the circle. \pracmarks{1 mark}
\end{enumerate}
\end{practiceq}

\begin{practiceq}{4}{5}\label{pq:raise}
A \qty{900}{kg} satellite is in a circular orbit \qty{350}{km} above the Earth's surface.
Its engine moves it to a new circular orbit with a period of \qty{4.00}{h}.
\begin{enumerate}[label=(\alph*)]
  \item Calculate the radius of the new orbit. \pracmarks{1 mark}
  \item Calculate the smallest energy needed to move the satellite to the new orbit.
        \pracmarks{2 marks}
  \item Calculate the change in the satellite's kinetic energy. State what this means for its
        speed. \pracmarks{1 mark}
\end{enumerate}
\end{practiceq}

% =====================================================================
\Needspace{14\baselineskip}
\section*{Section III: Extended response}

\begin{practiceq}{4}{3}\label{pq:workenergy}
A small ball is tied to a string. In situation A, it moves in a horizontal circle on a
smooth (frictionless) table. In situation B, it moves in a vertical circle. In both
situations, the other end of the string is tied to a fixed point at the centre of the circle,
and there is no air resistance.

Using the ideas of work and energy, explain why the speed of the ball stays constant in
situation A but changes in situation B.
\end{practiceq}

\begin{practiceq}{4}{4}\label{pq:leogeo}
Compare low Earth orbits and geostationary orbits. In your answer, refer to their altitude
and period, and link one use of each orbit to a property of that orbit.
\end{practiceq}

\begin{practiceq}{5}{2}\label{pq:stopper}
A student wants to show that the centripetal force on an object is proportional to the
square of its speed ($F_\mathrm{c} \propto v^2$). She ties a rubber stopper to a string. The
string passes through a short plastic tube held in her hand. A hanging mass $M$ on the other
end of the string gives the tension, $F_\mathrm{c} = Mg$. This is her method.
\begin{enumerate}
  \item Measure the mass of the stopper.
  \item Hang a \qty{100}{g} mass on the string. Whirl the stopper above your head.
  \item Use a stopwatch to time \emph{one} revolution of the stopper.
  \item Change the hanging mass to \qtylist{150;200;250;300}{g}. Time one revolution for each.
  \item After all the trials, measure the string from the top of the tube to the stopper. Use
        this length as the radius for every trial.
\end{enumerate}
Evaluate the validity and reliability of this method. Recommend improvements.
\end{practiceq}

% layout QA: start this question on a fresh page if it does not fit, so its text, table and
% grid stay together (it is about 44 lines tall)
\Needspace{44\baselineskip}
\begin{practiceq}{6}{1}\label{pq:launcher}
A student fixes a spring launcher on the floor. It fires a small ball at the same speed $u$
every time. The ball leaves the launcher at floor level. She changes the launch angle
$\theta$ and measures the range $R$ on the level floor. Each range in the table is the mean
of five trials.
\begin{table}
  \caption*{Range of the ball for each launch angle.}
  \centering
  \begin{tabular}{@{}lcccccc@{}}
    \toprule
    Launch angle $\theta$ & \ang{10} & \ang{15} & \ang{20} & \ang{25} & \ang{30} & \ang{45}\\
    \midrule
    $\sin 2\theta$ & \pracblank & \pracblank & \pracblank & \pracblank & \pracblank & \pracblank\\
    \midrule
    Mean range $R$ (\si{m}) & 0.660 & 0.958 & 1.236 & 1.482 & 1.661 & 1.918\\
    \bottomrule
  \end{tabular}
\end{table}
The model predicts $R = \dfrac{u^2\sin 2\theta}{g}$.
\begin{enumerate}[label=(\alph*)]
  \item Complete the table: calculate $\sin 2\theta$ for each angle (3 decimal places).
        \pracmarks{1 mark}
  \item On the grid below, plot $R$ against $\sin 2\theta$. Draw a line of best fit.
        \pracmarks{2 marks}
  \item Use the gradient of your line to calculate the launch speed $u$. \pracmarks{2 marks}
  \item The maker of the launcher says that $u = \qty{4.40}{m.s^{-1}}$. Suggest a reason for
        the difference between this value and your answer to (c). \pracmarks{1 mark}
\end{enumerate}
\begin{figure}
  \centering
  \begin{tikzpicture}
    \begin{axis}[hsc, width=0.86\linewidth, height=0.6\linewidth,
        xmin=0, xmax=1.1, ymin=0, ymax=2.2,
        xlabel={$\sin 2\theta$}, ylabel={$R$ (\si{m})},
        xtick={0,0.2,...,1.0}, ytick={0,0.4,...,2.0},
        minor x tick num=3, minor y tick num=3,
        grid=both, minor grid style={hscrule!45, line width=0.2pt},
        xticklabel style={/pgf/number format/.cd, fixed, fixed zerofill, precision=1},
        yticklabel style={/pgf/number format/.cd, fixed, fixed zerofill, precision=1}]
      \addplot[draw=none] coordinates {(0,0)};
    \end{axis}
  \end{tikzpicture}
  \caption*{Grid for part (b).}
\end{figure}
\end{practiceq}

% layout QA: keep the question text with its table and graph (about 39 lines tall)
\Needspace{39\baselineskip}
\begin{practiceq}{6}{4}\label{pq:trappist}
The star TRAPPIST-1 has seven known planets. Their orbits are almost circular.
The table gives data for five of them. The graph shows $r^3$ against $T^2$ for these five
planets.
\begin{table}
  \caption*{Orbits of five planets of TRAPPIST-1 ($r$ = orbital radius, $T$ = period).}
  \centering
  \begin{tabular}{@{}lcccc@{}}
    \toprule
    Planet & $r$ (\si{m}) & $T$ (\si{s}) & $r^3$ ($\times 10^{29}$ \si{m^3}) & $T^2$ ($\times 10^{12}$ \si{s^2})\\
    \midrule
    b & \num{1.73e9} & \num{1.31e5} & 0.0518 & 0.0172\\
    d & \num{3.33e9} & \num{3.50e5} & 0.369  & 0.123\\
    e & \num{4.38e9} & \num{5.27e5} & 0.840  & 0.278\\
    f & \num{5.76e9} & \num{7.96e5} & 1.91   & 0.634\\
    h & \num{9.26e9} & \num{1.62e6} & 7.94   & 2.62\\
    \bottomrule
  \end{tabular}
\end{table}
\begin{figure}
  \centering
  \begin{tikzpicture}
    \begin{axis}[hsc, width=0.78\linewidth, height=0.46\linewidth,
        xmin=0, xmax=3.0, ymin=0, ymax=9,
        xlabel={$T^2$ ($\times 10^{12}$ \si{s^2})}, ylabel={$r^3$ ($\times 10^{29}$ \si{m^3})},
        xtick={0,0.5,...,3.0}, ytick={0,1,...,9},
        minor x tick num=4, minor y tick num=1,
        grid=both, minor grid style={hscrule!45, line width=0.2pt},
        xticklabel style={/pgf/number format/.cd, fixed, fixed zerofill, precision=1}]
      \addplot[hscseriesB, line width=1pt, domain=0:2.9, samples=2, mark=none] {3.025*x};
      \addplot[hscseriesA, only marks, mark=*, mark size=2pt]
        coordinates {(0.0172,0.0518) (0.1225,0.3693) (0.2777,0.8403) (0.6336,1.911) (2.624,7.940)};
      \node[anchor=west, font=\small, text=hscink] at (axis cs:2.62,7.94) {\,h};
      \node[anchor=south east, font=\small, text=hscink] at (axis cs:0.6336,1.911) {f\,};
      \node[anchor=west, font=\small, text=hscseriesB] at (axis cs:1.55,3.3) {line of best fit};
    \end{axis}
  \end{tikzpicture}
  \caption*{$r^3$ against $T^2$ for five planets of TRAPPIST-1.}
\end{figure}
\begin{enumerate}[label=(\alph*)]
  \item Start from Newton's law of universal gravitation. Show that
        $\dfrac{r^3}{T^2} = \dfrac{GM}{4\pi^2}$ for a planet in a circular orbit around a star
        of mass $M$. \pracmarks{2 marks}
  \item Use the gradient of the line to calculate the mass of TRAPPIST-1.
        \pracmarks{3 marks}
  \item Planet g orbits TRAPPIST-1 at $r = \qty{7.01e9}{m}$. Predict its period in days.
        \pracmarks{1 mark}
\end{enumerate}
\end{practiceq}

% layout QA: keep all three parts of this question together (about 17 lines tall)
\Needspace{17\baselineskip}
\begin{practiceq}{7}{5}\label{pq:ellipse}
A \qty{400}{kg} satellite moves around the Earth in an elliptical orbit. At perigee (its
closest point) it is $r_\mathrm{p} = \qty{6.90e6}{m}$ from the Earth's centre, and its speed is
$v_\mathrm{p} = \qty{9.63e3}{m.s^{-1}}$. At apogee (its farthest point) it is
$r_\mathrm{a} = \qty{2.76e7}{m}$ from the Earth's centre.
\begin{enumerate}[label=(\alph*)]
  \item Use Kepler's second law to calculate the speed of the satellite at apogee.
        \pracmarks{1 mark}
  \item Calculate the total mechanical energy of the satellite at perigee and at apogee.
        Comment on your answers. \pracmarks{3 marks}
  \item Using forces and energy, explain why the satellite slows down as it moves from
        perigee to apogee. \pracmarks{3 marks}
\end{enumerate}
\end{practiceq}
